% Source: https://artofproblemsolving.com/community/c3h1225923_mandelbrot_competition_sample_national_level_test
% Stable identity: aops-topic-1225923-q01
$L_1$ and $L_2$ pass through the origin at angles $60^\circ$ and $30^\circ$. Reflecting in $L_2$ and then in $L_1$ is a rotation about the origin by $2(60^\circ-30^\circ)=60^\circ$. So $S$ is closed under rotations by multiples of $60^\circ$ and under reflection in $L_1$; these generate the symmetry group of the regular hexagon ($12$ elements: $6$ rotations and $6$ reflections in lines through $O$ at multiples of $30^\circ$).

The orbit of $(6,5)$ under this group has $12$ points unless $(6,5)$ lies on one of the six mirror lines. Its polar angle is $\arctan\dfrac{5}{6}\approx39.8^\circ$, not a multiple of $30^\circ$, so the orbit has $\boxed{12}$ points, and this orbit itself is a valid $S$.
\begin{asy}
usepackage("amsmath");
size(3.5cm);
defaultpen(fontsize(10pt));
draw((-1,-sqrt(3))--(1,sqrt(3)),blue); draw((-sqrt(3),-1)--(sqrt(3),1),red);
for(int k=0;k<6;++k){ dot(rotate(60k)*(1.2,1.0)); dot(rotate(60k)*reflect((0,0),(1,sqrt(3)))*(1.2,1.0),gray);}
label("$L_1$",(1,sqrt(3)),N); label("$L_2$",(sqrt(3),1),E); dot((1.2,1.0),red);
label("\mbox{}",(min(currentpicture,user=true)+max(currentpicture,user=true))/2,black); // keeps the text after the figure black when LaTeX typesets the labels inline
\end{asy}
